\documentclass[12pt,aspectratio=169,ignorenonframetext]{ltxbeamer}

\usepackage[spanish,es-nodecimaldot,es-tabla]{babel}

\title{Matemáticas\\Profesionales}
\subtitle{Curso de {\lmr\LaTeX}}
\author{Joar Buitrago\texorpdfstring{\thanks{jebuitragoc@unal.edu.co}}{ <jebuitragoc@unal.edu.co>}}
\date{}

\addbibresource{bibliography.bib}


\begin{document}
\begin{frame}
\titlepage
\end{frame}



\begin{frame}
    \frametitle{Bienvenidos al Nivel II}

    \centering
    {\fontsize{36}{38}\selectfont Pasar de documentos «bonitos» a documentos «científicamente precisos».\par}
\end{frame}



\section{El estándar de la AMS}
\begin{frame}
    \frametitle{¿Por qué no basta con el modo básico?}
    La \textit{American Mathematical Society} (AMS) desarrolló una serie de paquetes para resolver problemas que el {\lmr\TeX} original no manejaba de forma nativa:
    \vspace{1em}\pause
    \begin{itemize}
        \item<+-> \textbf{Alineación:} Ecuaciones que ocupan varias líneas.
        \item<+-> \textbf{Numeración:} Control fino sobre qué se numera y cómo.
        \item<+-> \textbf{Símbolos:} Acceso a la biblioteca completa de símbolos científicos con \texttt{amssymb}.
        \item<+-> \textbf{Estructura:} Entornos específicos para teoremas, lemas y pruebas.
    \end{itemize}
\end{frame}



\section{Alineación de Ecuaciones}
\subsection{align}
\begin{frame}[fragile]
\frametitle{El entorno \texttt{align}}
Permite alinear en columnas usando el símbolo \emph{et} (\texttt{\&}).

\begin{columns}
    \begin{column}{0.5\textwidth}
        \begin{block}{Código}
\begin{minted}{latex}
\begin{align}
  (a+b)^2 &= a^2 + 2ab + b^2 \\
  (a-b)^2 &= a^2 - 2ab + b^2
\end{align}
\end{minted}
        \end{block}
    \end{column}

    \begin{column}{0.5\textwidth}
        \begin{exampleblock}{Resultado}
            \begin{equation*}
                \begin{aligned}
                    (a+b)^2 &= a^2 + 2ab + b^2 \quad (1) \\
                    (a-b)^2 &= a^2 - 2ab + b^2 \quad (2)
                \end{aligned}
            \end{equation*}
        \end{exampleblock}
    \end{column}
\end{columns}
\bigskip
\small \textit{Nota: El símbolo \texttt{\&} define el punto de anclaje (usualmente el igual).}
\end{frame}

\begin{frame}[fragile]
\frametitle{El entorno \texttt{align}}
Permite alinear en columnas usando el símbolo \emph{et} (\texttt{\&}).

\begin{columns}
    \begin{column}{0.5\textwidth}
        \begin{block}{Código}
\begin{minted}{latex}
\begin{align*}
  \int \frac{2x}{x^2} \, dx
    &= \int \frac{1}{u} \, du
    & u &= x^2 \\
  &= \ln(u) + C \\
  &= \ln(x^2) + C
\end{align*}
\end{minted}
        \end{block}
    \end{column}

    \begin{column}{0.5\textwidth}
        \begin{exampleblock}{Resultado}
            \begin{equation*}
                \begin{aligned}
                    \int \frac{2x}{x^2} \, dx
                      &= \int \frac{1}{u} \, du
                      & u &= x^2 \\
                    &= \ln(u) + C \\
                    &= \ln(x^2) + C
                \end{aligned}
            \end{equation*}
        \end{exampleblock}
    \end{column}
\end{columns}
\end{frame}



\subsection{gather}
\begin{frame}[fragile]
\frametitle{El entorno \texttt{gather}}
Centra varias ecuaciones sin alinearlas entre sí.

\begin{columns}
\begin{column}{0.5\textwidth}
    \begin{block}{Código}
\begin{minted}{latex}
\begin{gather}
  a^2 \pm 2ab + b^2 = (a \pm b)^2  \\
  a^2 - b^2 = (a + b)(a - b)
\end{gather}
\end{minted}
        \end{block}
    \end{column}

    \begin{column}{0.5\textwidth}
        \begin{exampleblock}{Resultado}
            \begin{equation*}
                \begin{gathered}
                    a^2 \pm 2ab + b^2 = (a \pm b)^2 \quad (1) \\
                    a^2 - b^2 = (a + b)(a - b) \quad (2)
                \end{gathered}
            \end{equation*}
        \end{exampleblock}
    \end{column}
\end{columns}
\end{frame}

\begin{frame}[fragile]
\frametitle{El entorno \texttt{gather}}
Centra varias ecuaciones sin alinearlas entre sí.

\begin{columns}
    \begin{column}{0.5\textwidth}
        \begin{block}{Código}
\begin{minted}{latex}
\begin{gather*}
  \int \frac{2x}{x^2} \, dx
    = \int \frac{1}{u} \, du
  u = x^2 \\
  = \ln(u) + C \\
  = \ln(x^2) + C
\end{gather*}
\end{minted}
        \end{block}
    \end{column}

    \begin{column}{0.5\textwidth}
        \begin{exampleblock}{Resultado}
            \begin{equation*}
                \begin{gathered}
                    \int \frac{2x}{x^2} \, dx
                      = \int \frac{1}{u} \, du \\
                    u = x^2 \\
                    = \ln(u) + C \\
                    = \ln(x^2) + C
                \end{gathered}
            \end{equation*}
        \end{exampleblock}
    \end{column}
\end{columns}
\end{frame}



\subsection{multline}
\begin{frame}[fragile]
\frametitle{El entorno \texttt{multline}}
Para una sola ecuación que es demasiado larga para una línea.

\begin{columns}
\begin{column}{0.5\textwidth}
    \begin{block}{Código}
\begin{minted}{latex}
\begin{multline}
  A = a + b + c + d + e \\
  + f + g + h + i + j + k \\
  + l + m + n + o + p
\end{multline}
\end{minted}
        \end{block}
    \end{column}

    \begin{column}{0.5\textwidth}
        \begin{exampleblock}{Resultado}
            \lmr
            \begin{multline*}
                A = a + b + c + d + e \\
                + f + g + h + i + j + k \\
                + l + m + n + o + p
            \end{multline*}\\[-1em]
            \hfill(1)
        \end{exampleblock}
    \end{column}
\end{columns}
\end{frame}

\begin{frame}[fragile]
\frametitle{El entorno \texttt{multline}}
Para una sola ecuación que es demasiado larga para una línea.

\begin{columns}
\begin{column}{0.5\textwidth}
    \begin{block}{Código}
\begin{minted}{latex}
\begin{multline*}
  A = a + b + c + d + e \\
  + f + g + h + i + j + k \\
  + l + m + n + o + p
\end{multline*}
\end{minted}
        \end{block}
    \end{column}

    \begin{column}{0.5\textwidth}
        \begin{exampleblock}{Resultado}
            \begin{multline*}
                A = a + b + c + d + e \\
                + f + g + h + i + j + k \\
                + l + m + n + o + p
            \end{multline*}
        \end{exampleblock}
    \end{column}
\end{columns}
\end{frame}



\subsection{aligned, gathered, split}
\begin{frame}[fragile]
\frametitle{Ecuaciones numeradas como bloque}
A veces queremos una sola numeración para varias líneas de desarrollo.

La \emph{AMS} nos ofrece los entornos:

\begin{itemize}
    \item \texttt{aligned} \uncover<2->{\(\rightarrow\) \texttt{align}}
    \item \texttt{gathered} \uncover<2->{\(\rightarrow\) \texttt{gather}}
    \item \texttt{split} \uncover<2->{\(\rightarrow\) \texttt{align}\alert{*}}
\end{itemize}

\pause
Que funcionan igual que su contraparte, pero solamente se pueden utilizar dentro de otro entorno.
\end{frame}

\begin{frame}[fragile]
\frametitle{Ecuaciones numeradas como bloque}

\begin{columns}
    \begin{column}{0.5\textwidth}
        \begin{block}{Código}
\begin{minted}{latex}
\begin{equation}
  \begin{aligned}
    a + b &= 7 \\
    ab &= 10
  \end{aligned}
\end{equation}
\end{minted}
        \end{block}
    \end{column}

    \begin{column}{0.5\textwidth}
        \begin{exampleblock}{Resultado}
            \begin{equation*}
                \begin{aligned}
                    a + b &= 7 \\
                    ab &= 10
                \end{aligned}
                \quad (1)
            \end{equation*}
        \end{exampleblock}
    \end{column}
\end{columns}
\end{frame}



\section{Matrices y Arreglos}
\begin{frame}
\frametitle{Matrices de la AMS}

Los paquetes de la \emph{AMS} facilitan la creación de matrices con diferentes delimitadores:
\begin{itemize}
    \item \texttt{matrix}: sin nada
    \item \texttt{pmatrix}: paréntesis
    \item \texttt{bmatrix}: corchetes
    \item ...
\end{itemize}
\end{frame}

\begin{frame}[fragile]
\frametitle{Matrices de la AMS}

\begin{columns}
    \begin{column}{0.5\textwidth}
        \begin{block}{Código}
\begin{minted}{latex}
\begin{pmatrix}
  1 & 0 & 0 \\
  0 & 1 & 0 \\
  0 & 0 & 1
\end{pmatrix}
\end{minted}
        \end{block}
    \end{column}

    \begin{column}{0.5\textwidth}
        \begin{exampleblock}{Resultado}
            \begin{equation*}
                \begin{pmatrix}
                    1 & 0 & 0 \\
                    0 & 1 & 0 \\
                    0 & 0 & 1
                \end{pmatrix}
            \end{equation*}
        \end{exampleblock}
    \end{column}
\end{columns}
\end{frame}



\begin{frame}[fragile]
\frametitle{Matrices de la AMS}

\begin{columns}
    \begin{column}{0.5\textwidth}
        \begin{block}{Código}
\begin{minted}{latex}
\begin{bmatrix}
  1 & 0 \\
  0 & 1 \\
  0 & 0
\end{bmatrix}
\end{minted}
        \end{block}
    \end{column}

    \begin{column}{0.5\textwidth}
        \begin{exampleblock}{Resultado}
            \begin{equation*}
                \begin{bmatrix}
                    1 & 0 \\
                    0 & 1 \\
                    0 & 0
                \end{bmatrix}
            \end{equation*}
        \end{exampleblock}
    \end{column}
\end{columns}
\end{frame}



\section{Definiciones por casos}
\begin{frame}[fragile]
    \frametitle{Definiciones por casos}
    Fundamental para funciones definidas a trozos.

    \begin{columns}
        \column{0.5\textwidth}
\begin{minted}{latex}
f(x) = \begin{cases}
  1 & \text{si } x \in \mathbb{Q} \\
  0 & \text{si \(x \notin \mathbb{Q}\)}
\end{cases}
\end{minted}

        \column{0.5\textwidth}
        \begin{exampleblock}{Resultado}
            \begin{equation*}
                f(x) = \begin{cases}
                    1 & \text{si } x \in \mathbb{Q} \\
                    0 & \text{si  \(x \notin \mathbb{Q}\)}
                \end{cases}
            \end{equation*}
        \end{exampleblock}
    \end{columns}
\end{frame}



\begin{frame}
    \frametitle{Laboratorio de hoy}
    Replicar el desarrollo de una identidad trigonométrica y una serie de Taylor usando los entornos \texttt{align} y \texttt{split}.

    \bigskip
    \textbf{Requisitos:}
    \begin{itemize}
        \item El documento debe usar el paquete \texttt{amsmath}.
        \item Debe haber al menos una ecuación con numeración personalizada.
        \item Debe incluirse una matriz de transformación.
    \end{itemize}
\end{frame}



\begin{frame}[allowframebreaks]
\frametitle{Referencias}
\nocite{*}
\printbibliography
\end{frame}
\end{document}
